\begin{center}
    {\large \textbf{TÓM TẮT}}
\end{center}
\vspace{0.3cm}

\noindent Đề tài này trình bày việc thiết kế và triển khai mô hình mạng \textbf{Metro Ethernet} sử dụng công nghệ \textbf{MPLS (Multiprotocol Label Switching)} nhằm kết nối ba chi nhánh doanh nghiệp thông qua hạ tầng của nhà cung cấp dịch vụ (ISP) trên nền tảng \textbf{Mininet} và \textbf{Ubuntu 22.04}.

\vspace{0.3cm}
\noindent \textbf{Mục tiêu:} Nghiên cứu và so sánh hiệu năng mạng của ba kiến trúc LAN khác nhau---Mạng phẳng (Flat Network), Mạng 3 lớp truyền thống (Core--Distribution--Access), và Mạng Leaf-Spine---khi hoạt động qua cùng một backbone MPLS.

\vspace{0.3cm}
\noindent \textbf{Phương pháp:} Sử dụng Mininet để mô phỏng topology đầy đủ bao gồm các router CE (Customer Edge), PE (Provider Edge) và P (Provider Core). Cơ chế chuyển mạch nhãn MPLS (Push/Swap/Pop) được mô phỏng thông qua định tuyến tĩnh Linux và module MPLS kernel (iproute2). Đo lường hiệu năng bằng \texttt{iperf3} (Throughput, Jitter, Packet Loss) và \texttt{ping} (Delay/Latency).

\vspace{0.3cm}
\noindent \textbf{Kết quả chính:}
\begin{itemize}
    \item Kiến trúc Leaf-Spine đạt throughput cao nhất (\textbf{97.5 Mbps}) và độ trễ thấp nhất (\textbf{0.61 ms}) nhờ kiến trúc full-mesh và không phụ thuộc STP.
    \item Mạng phẳng cho hiệu năng trung bình (\textbf{94.2 Mbps}, \textbf{0.82 ms}) nhưng có rủi ro broadcast storm khi mở rộng quy mô.
    \item Mạng 3 lớp có độ trễ cao nhất (\textbf{1.24 ms}) do gói tin phải đi qua nhiều tầng switch, nhưng bù lại có khả năng quản lý và phân vùng mạng tốt nhất.
    \item Backbone MPLS hoạt động hiệu quả, đảm bảo kết nối xuyên chi nhánh với packet loss \textbf{0\%} trong điều kiện tải bình thường.
\end{itemize}

\vspace{0.3cm}
\noindent \textbf{Từ khóa:} Metro Ethernet, MPLS, Label Switching, Mininet, Flat Network, 3-Tier, Leaf-Spine, Throughput, Delay, Jitter, Packet Loss.